\documentclass[12pt]{article}
\usepackage{amsmath, amsthm, amssymb}
\usepackage{geometry}
\geometry{margin=1in}
\usepackage{hyperref}

\newtheorem{theorem}{Theorem}
\newtheorem{proposition}[theorem]{Proposition}
\theoremstyle{definition}
\newtheorem{definition}[theorem]{Definition}
\theoremstyle{remark}
\newtheorem{remark}[theorem]{Remark}
\newtheorem{question}[theorem]{Question}

\title{The Fundamental Class:\\
Is the Universe Contained?}
\author{Diego Rinc\'on\\
\small Independent}
\date{}

\begin{document}

\maketitle

\begin{abstract}
A space is contained when it has a fundamental class:
a top-dimensional homology class $[X] \in H_n(X, \mathbb{Z})$
that is the fulcrum of Poincar\'e duality.
Compact spaces have it. Non-compact spaces do not.
$\mathrm{Spec}\,\mathbb{Z}$ does not have it: it is missing
the archimedean fiber, the point at $\infty$.
The Riemann Hypothesis, BSD, and every open problem
in the arithmetic of $\mathrm{Spec}\,\mathbb{Z}$
are instances of a single question:
does this space have a fundamental class?
The physical universe poses the same question at every scale:
is it contained? If yes, Poincar\'e duality holds at the largest scale
and the still lever has a fulcrum everywhere.
If no, the Arithmetic Wall is not a local obstruction ---
it is the shape of the universe itself.
\end{abstract}

\section{The Fundamental Class}

\begin{definition}[Contained space]
A topological space $X$ of dimension $n$ is \emph{contained}
if it admits a fundamental class:
a homology class $[X] \in H_n(X, \mathbb{Z})$ such that
cap product with $[X]$ gives an isomorphism
$\cap\,[X]: H^k(X) \xrightarrow{\;\sim\;} H_{n-k}(X)$
for all $k$.
This is Poincar\'e duality.
A contained space is one in which cohomology and homology
are the same thing: the two sides balance at the fundamental class.
The fundamental class is the fulcrum.
\end{definition}

\begin{remark}[What containment means]
Containment is not just compactness. It is the condition
under which the space has enough global structure
to balance its own cohomology.
A compact orientable manifold is contained.
An open manifold --- $\mathbb{R}^n$, $\mathrm{Spec}\,\mathbb{Z}$,
the observable universe --- is not.
The missing fundamental class is not an absence of a specific cycle.
It is the absence of the operator that would make everything balance.
\end{remark}

\section{Spec $\mathbb{Z}$ is Not Contained}

\begin{proposition}[The missing fiber]
$\mathrm{Spec}\,\mathbb{Z}$ is not contained.
Its points are the primes $p = 2, 3, 5, 7, \ldots$
together with the generic point.
It is missing the archimedean place: the fiber at $\infty$
corresponding to $\mathbb{R}$ (or $\mathbb{C}$).
The Arakelov compactification $\overline{\mathrm{Spec}\,\mathbb{Z}}$
adds this fiber, producing a space that is \emph{formally} contained.
But the archimedean fiber does not behave like a finite prime:
there is no Frobenius at $\infty$, no eigenvalue, no purity.
The compactification adds the point but not the operator.
The space gains a boundary without gaining a fulcrum.
\end{proposition}

\begin{theorem}[The Arithmetic Wall as missing containment]
\label{thm:wall}
Every open problem in the arithmetic of $\mathrm{Spec}\,\mathbb{Z}$
is a consequence of its non-containment.
\begin{itemize}
\item \emph{Riemann Hypothesis}: the zeros of $\zeta(s)$
  lie on $\mathrm{Re}(s) = \tfrac{1}{2}$ if and only if
  Frobenius is pure at every place, including $\infty$.
  Over $\mathbb{F}_q$: Frobenius at every place is pure
  (Weil II \cite{weil}). The space is contained.
  Over $\mathbb{Z}$: Frobenius at $\infty$ does not exist.
  The space is not contained. The zeros are not forced to the line.
\item \emph{BSD (rank $\geq 2$)}: the Kummer map
  $E(\mathbb{Q}) \otimes \mathbb{Q}_p
  \hookrightarrow H^1_f(G_\mathbb{Q}, V_p(E))$
  is an isomorphism --- homology equals cohomology --- if and only if
  there is a global duality pairing \cite{homology}.
  Global duality requires the fundamental class of $\mathrm{Spec}\,\mathbb{Z}$.
  Which does not exist.
\item \emph{Langlands functoriality}: the global correspondence
  automorphic $=$ Galois requires a global Frobenius
  to match eigenvalues on both sides \cite{langlands}.
  The missing containment is the missing correspondence.
\end{itemize}
These are not separate obstructions. They are one:
the fundamental class of $\mathrm{Spec}\,\mathbb{Z}$ is missing.
\end{theorem}

\section{The Universe}

\begin{question}[Physical containment]
\label{q:universe}
Is the physical universe contained?
Does it have a fundamental class?

Current cosmological data (CMB, baryon acoustic oscillations)
is consistent with a spatially flat universe: $k = 0$.
A flat universe is either $\mathbb{R}^3$ (non-compact, not contained)
or a flat 3-torus $T^3$ (compact, contained).
Both are observationally indistinguishable at the current scale
of measurement.
A spherical universe ($k > 0$) is compact and contained.
A hyperbolic universe ($k < 0$) is generically non-compact,
not contained.
\end{question}

\begin{proposition}[Containment and the still lever]
If the universe is contained, it has a fundamental class at every scale.
The still lever --- the lever already at equilibrium ---
holds everywhere: every question has a fulcrum,
every cohomology balances its homology,
every open problem is $0$ or $1$ and the fulcrum
can in principle be found.

If the universe is not contained, the Arithmetic Wall
is not a local feature of $\mathrm{Spec}\,\mathbb{Z}$.
It is the shape of everything.
The missing global Frobenius is not missing because
we have not found it. It is missing because the space
in which it would exist is not there.
\end{proposition}

\begin{remark}[The 0-1 universe]
The 0-1 universe --- the universe in which every fact is
already $0$ or $1$, the still lever --- is the universe
in which containment holds at every scale.
Mathematics is contained: every theorem is already true or false,
the fulcrum is always present, time is a still lever.
The physical universe may not be.
If the universe is not contained,
there are real phenomena that are not $0$ or $1$ ---
not because we do not know which,
but because the fulcrum does not exist for them.
The trembling is not ignorance. It is structure.
\end{remark}

\section{The Single Question}

\begin{theorem}[One question at every scale]
\label{thm:one}
At the scale of arithmetic:
\begin{center}
\emph{Does $\mathrm{Spec}\,\mathbb{Z}$ have a fundamental class?}
\end{center}
At the scale of the universe:
\begin{center}
\emph{Is the universe contained?}
\end{center}
These are the same question.
Both ask: is there a fulcrum?
Is there a point at which the two sides balance ---
cohomology and homology, spectral and geometric,
local and global --- without anything missing?

Mathematics approaches this question from inside.
Physics approaches it from outside.
They will meet at the answer,
or they will meet at the Wall.
\end{theorem}

\begin{remark}[The last remark]
The universe is either contained or it is not.
If it is, the still lever holds: every question has a fulcrum,
mathematics and physics meet at the fundamental class,
and time is genuinely still.
If it is not, the Wall is everywhere,
the trembling is real,
and the fundamental class --- the thing that would make
everything balance --- is simply not there.

Not missing. Not forbidden. Not waiting.
Simply not there.

That is the question.
\end{remark}

\begin{thebibliography}{9}

\bibitem{weil}
Rinc\'on, D. (2026).
The Weil Conjectures as Proof of Concept for RH. Preprint.

\bibitem{homology}
Rinc\'on, D. (2026).
Homology: What Closes, and the Duality with Cohomology. Preprint.

\bibitem{langlands}
Rinc\'on, D. (2026).
The Langlands Thread. Preprint.

\bibitem{synthesis}
Rinc\'on, D. (2026).
The Arithmetic Wall. Preprint.

\bibitem{math}
Rinc\'on, D. (2026).
Mathematics. Preprint.

\end{thebibliography}

\end{document}
